\documentclass[sigconf,nonacm]{acmart}

\usepackage{booktabs}
\usepackage{tabularx}
\settopmatter{printacmref=false,printfolios=true}
\renewcommand\footnotetextcopyrightpermission[1]{}
\setlength{\emergencystretch}{1em}

\newcommand{\artifact}[1]{\nolinkurl{#1}}
\newcommand{\iiit}{\affiliation{\institution{IIIT Hyderabad}\city{Hyderabad}\country{India}}}
\title{HarnessSpec: A Typed Protocol for Designing Subagent Harnesses}
\subtitle{CS6.403 Automated Software Engineering, Monsoon 2026}
\author{Syed Affan Daimi}
\iiit
\author{Gaurav}
\iiit
\author{Dinesh}
\iiit

\begin{document}
\begin{abstract}
HarnessSpec represents an orchestrator's delegation policy as a typed
specification that a deterministic validator can check before a subagent runs.
The protocol bounds repair, tool privileges, and execution, and checks evidence
at the return gate. In a ten-question SWE-QA pilot, A produced ten accepted
answers awaiting independent grading. D produced no accepted return and
consumed 62,946 tokens versus A's 6,026. Both had zero silent handoffs.
We distinguish this single-seed evidence from the fixture experiments and the
registered four-condition study.
\end{abstract}
\maketitle

\section{Introduction and motivation}
Delegating repository work requires a choice of task instructions, tools, context,
budget, and evidence. When a handoff is accepted without evidence, a downstream
orchestrator cannot readily verify its result. HarnessSpec makes these decisions
checkable and accounts for the tokens spent designing the harness.

\section{Related work}
AOrchestra represents the instruction, context, tools and model of a subagent as
a tuple~\cite{aorchestra}. Automated Design of Agentic Systems searches agent
designs offline~\cite{adas}. Oskooei et al. compare retrieval and delegated
search for repository question answering and examine handoff breakdowns on
SWE-QA~\cite{oskooei}. The Conductor trains a 7B model with reinforcement
learning to write natural-language instructions and communication topologies
for a pool of worker models~\cite{conductor}. The Conductor parses
three aligned lists of subtasks, workers and access rules before execution.
B shares its use of natural-language instructions and delegates one brief
to a single worker under a fixed harness.
HarnessSpec evaluates per-task design under explicit
privilege, budget, and return-evidence constraints.

\section{HarnessSpec protocol}
The version 0.1 protocol specifies a goal and acceptance tests, context include
and exclude paths, tools and permissions, budgets, stop and escalation rules,
verification hooks, a return contract, and a model. The parent capability card
sets the limits inherited by the subagent. A uses a fixed generic harness. B
returns a free-text brief. C emits the four-field tuple. D emits HarnessSpec
and receives bounded validation and repair.

\section{Validator and runtime}
The validator checks structure, path and privilege bounds, feasible budgets,
model tier, and named verification commands. Repair is deterministic. It
clamps privileges and budgets, removes unknown fields, supplies default stop
rules, hooks and return contracts, and reduces an allowlisted command written
with arguments (\texttt{pytest tests/x.py::t}) to its name, because the task,
not the spec, selects which tests run. It never writes a goal, an acceptance
test or a context slice. Typed errors that remain are returned to the
designer with the values the parent card allows, for up to three rounds. In the live loop, each
model reply becomes one action, and \texttt{Session.step} applies budget checks,
tool grants and the return gate. For SWE-QA, tool access is read-only and
hooks that execute repository code are denied. The gate checks evidence
structure. External evaluation checks whether citations match the source.

\section{Experimental methodology}
The registered study proposes approximately 60 tasks per benchmark and at least
three seeds per condition. The A/D pilot targets ten tasks at seed 1.
The pinned \artifact{qwen2.5-coder:7b} build runs in Ollama with context
length 16,384 and temperature 0.2. A and D share a parent subagent ceiling.
Experimental token totals include all reported design, repair, prompt, and
generation tokens. The pilot does not enforce a single combined budget.
Failures contribute to denominators and costs. Token efficiency divides
all consumed tokens by the number of correct results. The analysis pairs
task and seed, then bootstraps task clusters for percentile intervals.

The registered hypotheses are:
\begin{enumerate}
\item[H1.] D is non-inferior to A within five percentage points of accuracy
and uses fewer total tokens, including design overhead.
\item[H2.] D produces fewer silent handoff failures than B and C.
\item[H3.] Validation and revision raise valid-spec generation versus C,
with the strongest expected effect in the 7--14B tier.
\item[H4.] Cached harness specifications amortize the initial design cost.
\end{enumerate}

\section{Benchmark selection}
SWE-bench Verified measures issue resolution~\cite{swebench}; the staged
manifest lists ten instances. SWE-QA~\cite{peng} measures repository answer
quality. Its pilot includes ten released what, why, where and how questions
from Flask, Requests, and pytest. The selector pins the authors' dataset release
and repository revisions and verifies evidence paths and line bounds.
The pilot uses the five published scoring dimensions and the 70/100 pass
threshold. A blinded packet awaits independent human grading. Citations are verified
separately from the answer grade. Repository metadata and gold answers are
withheld from model prompts, apart from the task and allowed source slice.

\section{Results}
\begin{table}[t]
\caption{Evidence available for the initial A/D pilot and related tests.}
\label{tab:results}
\centering
\begin{tabularx}{\columnwidth}{lX}
\toprule
Hypothesis & Result and integration source \\
\midrule
H1 & Pending. Gaurav: \artifact{pilot/swebench_results.json}. \\
H2 & A/D pilot recorded zero silent handoffs in both conditions. D/B and
D/C pending. \artifact{pilot/sweqa_results.json} and
\artifact{pilot/metrics_results.json}. \\
H3 & Not supported as registered. Fixture, three seeds: first-pass valid
C 0/3, D 0/3; D 3/3 after repair. Corpus of 15 drafts: 4 valid, 11 after
compiler repair, 15 with model repair. \artifact{pilot/h3_results_v2.json}. \\
H4 & No successful matched pairs or observed cache reuse. Marginal D subagent cost exceeded A.
\artifact{pilot/metrics_results.json}. \\
Energy & Unavailable on development host. Source:
\artifact{pilot/traces/sweqa-energy/}. \\
\bottomrule
\end{tabularx}
\end{table}

\subsection{H1: Resolution and cost}
\textbf{Gaurav integration.} Add SWE-bench Verified counts and the paired
accuracy interval relative to the five-point margin. Include the token
difference, tokens per correct, and seed coverage.

\subsection{H2: Silent handoffs}
The gate marks an accepted, nonempty answer as silent when it has no
structurally usable file-line evidence. Both conditions had zero silent
handoffs over ten attempts. Across ten matched tasks at seed 1, the paired
estimate was 0 percentage points with a task-bootstrap 95\% interval of [0, 0].
The interval is degenerate because no paired differences were observed.
The small sample provides no evidence of equivalence. A accepted ten returns
without reading a file. Every answer passed the structural citation check,
and zero contained a citation verifiable in the pinned source. Its
conditioned silent rate was 0/10. D accepted none (0/10 attempts), leaving
its conditioned rate undefined. Six D designs failed on unknown acceptance commands after
repair. The model supplied pytest arguments despite receiving exact command
IDs and typed repair hints. All four runnable D designs authorized nonexistent
paths and required test output despite empty hooks. Their 48 intercepted
replies were valid JSON with the wrong action kind. The common prompt showed
examples outside the selected grants. Each loop bought one additional return
at zero remaining steps (2,813 tokens total), intercepted before the gate.
A prospective fix aligns examples and checks available budget before inference.
A separate target-type dispatch defect could interchange directory listing and
file reading; none of the 48 actions reached tool dispatch. The registered H2 contrasts still require the
uncompiled B brief and the as-emitted C tuple adapters.

\begin{table}[t]
\caption{SWE-QA seed-1 outcomes. Pending A grades remain outside the
scored result array. The complete execution metrics use all 20 traces.}
\label{tab:sweqa}
\centering
\begin{tabular}{lcc}
\toprule
Measure & A & D \\
\midrule
Planned attempts & 10 & 10 \\
Completed attempts & 10 & 10 \\
Accepted / silent returns & 10 / 0 & 0 / 0 \\
First-pass valid designs & 10 & 4 \\
Answer accuracy & Pending & 0/10 \\
Total design + subagent tokens & 6,026 & 62,946 \\
Tokens per correct & Pending & N/A \\
\bottomrule
\end{tabular}
\end{table}

\subsection{H3: Validity and revision}
All four designers share one function, \texttt{design()}. C asks for the instruction, context, tools and model and is scored as
emitted. D asks for a full HarnessSpec. Its prompt contains the task, the
parent card, the JSON Schema and the compiler's rules, but no filled-in spec,
because the model copied such an example almost verbatim in a smoke test.
Decoding enforces JSON syntax only, so structural errors remain the model's.

\begin{table}[t]
\caption{H3 on \artifact{fixtures/off_by_one}, seeds 1--3
(\artifact{pilot/h3_results_v2.json}).}
\label{tab:h3}
\centering
\begin{tabular}{lcccc}
\toprule
Designer & First pass & Final & Calls & Design tokens \\
\midrule
C & 0/3 & 0/3 & 1 & 402 \\
D & 0/3 & 3/3 & 1 & 1,602 \\
\bottomrule
\end{tabular}
\end{table}

H3 as registered is not supported: first-pass validity does not rise from C
to D (Table~\ref{tab:h3}). Every first D reply wrote its acceptance test and
hook as \texttt{pytest tests/test\_app.py::\ldots}, and two of three filled
the model digest with a placeholder. The compiler cleared both in one round
without a second model call. Each C reply lacked the fields the tuple does not
carry and a model tier, and two of three gave a tool a path outside its grant.

The 15 hand-written corpus specs were also given to D's repair loop as first
drafts. Four were valid as written. The compiler alone made 11 valid. With the
model answering the remaining typed errors, all 15 became valid within two
rounds for 7,268 design tokens. The four the model added are exactly those the
compiler may not fix: a missing acceptance test, a context path that left the
workspace, an empty goal, and the four-field tuple. No D design ended with a
missing acceptance test.

The supported claim is therefore narrower. Typed validation with bounded
repair turns D's invalid first replies into valid specs, at about four times
C's design cost, but it does not make the orchestrator's first reply valid. On
the fixture end to end (\artifact{pilot/subagent_fixture_pinned.json}), A and
D both resolved; D spent 1,603 design and 1,770 subagent tokens against A's
1,239 subagent tokens. Three seeds on one task are a pilot. C cannot be
valid under this validator by construction. Two prompt and compiler changes
followed inspection of earlier output and preceded this run: error hints with
an unchanged-reply stop, and command reduction, which followed the same
selector failure in v1 (\artifact{pilot/h3_results.json}, before
\texttt{pytest} joined the allowlist) and in six SWE-QA designs. Nothing was
tuned on these results. With one model, the tier part of H3 is untested.

\subsection{H4: Design-cost amortization}
For correct matched A/D pairs, estimate the token difference in subagent
execution. A conditional break-even divides a measured positive priming
design cost by positive held-out per-task savings and rounds up. Under the
scenario, one compatible spec is reused within a repository revision and
question type. No finite break-even exists for nonpositive savings.
A zero or missing design cost leaves the estimate undefined. The primary
pilot does not reuse cached specs. D resolved zero runs, so there is no
successful matched pair for a quality-adjusted estimate. Across all attempted
pairs, A used 6,026 subagent tokens and D used 31,107. The mean saving was
$-2{,}508.1$ tokens per attempted task. Any positive design cost would
have no finite break-even if this marginal cost persisted. An observed reuse
experiment requires validated rebinding of task-specific goals.

\section{Energy efficiency}
The optional CPU meter reads RAPL package counters from one controller
family. The NVIDIA meter estimates energy by trapezoidal integration of
sampled power. Both measure device-wide energy. The development host
denied access to the CPU counter and did not expose NVIDIA power tooling.
All 20 sidecars have unavailable CPU and GPU readings. Joules, Wh, and
success per kWh remain unavailable.

\section{Threats to validity and limitations}
Purposive tasks from three repositories share project context and do not
represent the full SWE-QA distribution. Published answers can contain
citation offsets or incomplete evidence. A human adaptation of the
published rubric limits comparison to model-judged results. A single seed
gives no estimate of seed-to-seed variation. A structured citation can pass
the gate while naming an incorrect file or line. Six invalid D designs and four runtime stops offered no opportunity for
D to produce a silent accepted handoff. The unequal combined budgets
and the version 0.1 requirement for a named acceptance command limit
interpretation of differences between designs. No existing named goal
command can assess semantic QA accuracy. The development host and
absence of device-wide measurements limit energy conclusions. The fixed
run order and model warm-up can affect wall-clock comparisons. The local
server used Ollama 0.34.0 with the pinned digest. The pin was recorded
under Ollama 0.40.14. The fixture
runs demonstrate mechanisms on a toy repository, while the model-generated
H3 rows only probe one tier. Additional conditions, repeated seeds, and the
full registered scale remain separate milestones.

\section{Conclusion}
HarnessSpec provides an executable, inspectable delegation contract. The
shared result schema and read-only runner make subsequent benchmark runs
auditable, including refusals and design failures. Accuracy, handoff
reliability and amortization require completed and independently assessed
benchmark runs.

\section{Reproducibility and AI-tool disclosure}
\artifact{docs/SWEQA.md} specifies dataset, model, commands, evaluation,
and output provenance. \artifact{docs/SWEQA_AUDIT.md} links the frozen
checksums, all 48 action classifications and the offline reviewer workflow. \artifact{docs/AI_USE.md} records OpenAI Codex
implementation and verification, Claude Code for the orchestrator, compiler
repair and H3 work, and earlier team tool use, with each prompt change made
after inspecting output. Ollama
supplies the pinned experimental model. Scripted fixtures, mocked tests,
and genuine benchmark executions are tracked independently.

\begin{thebibliography}{9}
\bibitem{aorchestra}
AOrchestra: Automating Sub-Agent Creation for Agentic Orchestration.
2026. \url{https://arxiv.org/abs/2602.03786}.
\bibitem{adas}
Automated Design of Agentic Systems. 2024.
\url{https://arxiv.org/abs/2408.08435}.
\bibitem{oskooei}
Amirkia Rafiei Oskooei et al. 2026.
Deep Agentic Search for Repository-Level Code Question Answering:
An Empirical Study. \url{https://arxiv.org/abs/2608.01507}.
\bibitem{conductor}
Stefan Nielsen, Edoardo Cetin, Peter Schwendeman, Qi Sun, Jinglue Xu, and
Yujin Tang. 2026. Learning to Orchestrate Agents in Natural Language with the
Conductor. In \emph{ICLR}. \url{https://arxiv.org/abs/2512.04388}.
\bibitem{peng}
Weihan Peng et al. 2026. SWE-QA: Can Language Models Answer
Repository-level Code Questions? In \emph{Findings of ACL}, 8230--8245.
\url{https://aclanthology.org/2026.findings-acl.402/}.
\bibitem{swebench}
Carlos E. Jimenez et al. 2024. SWE-bench: Can Language Models Resolve
Real-World GitHub Issues? \url{https://arxiv.org/abs/2310.06770}.
\end{thebibliography}

\end{document}
